\begin{problem}
  Let $\mathcal{A} = \{a \in \mathds{R}^2 : \|\|a\|\|_2 \leq1 \}$,
  but let the $\infty$-norm be used as distance function.

\begin{enumerate}
\item Draw a diagram to show the best approximation in $\mathcal{A}$
  to a general point in $\mathds{R}^2$.  instead of a number.
\item Verify that the best approximation operator from $\mathds{R}^n$
  to $\mathcal{A}$ is continuous.
\end{enumerate}

\end{problem}

\begin{solution}
\begin{enumerate}
\item The best approximation to a general point $f \in \mathds{R}^2$
  the intersection between the line $\alpha(f) + x$ that go though $f$
  and the unit circle. Such an intersection does not exist for all
  points. If that is the case $f$ will be mapped to $(1, 0), (0, 1),
  (-1, 0), (0, -1)$ if it is in the top, right, bottom, or left side
  of the plane respectively.
  \begin{figure}[!ht]
  \centering \includegraphics[scale = 0.1]{drawing_task_4.png}
  \label{fig:task_4}
  \end{figure}

\item As the function $\alpha(f)$ is continuous so will also the
  intersection point and thus the best approximation by this
  mechanism. Further points not having such and intersection will all
  get the same best approximation as their neighbors, thus satisfying
  continuity. Last but not least on the border between the two methods
  of
\end{enumerate}
\end{solution}

 	

%%% Local Variables:
%%% TeX-master: "report.tex"
%%% End:
